\documentclass{ifpr-pinhais}

\usepackage{booktabs}
\usepackage{listings}
\usepackage{xcolor}
\usepackage{amsmath}
\usepackage{float}

\graphicspath{{img/}}

\definecolor{codebg}{HTML}{F4F4F4}
\definecolor{kw}{HTML}{00008B}
\definecolor{cm}{HTML}{006400}
\definecolor{str}{HTML}{8B0000}

\lstset{
    backgroundcolor=\color{codebg},
    frame=single,
    basicstyle=\ttfamily\footnotesize,
    keywordstyle=\color{kw}\bfseries,
    commentstyle=\color{cm}\itshape,
    stringstyle=\color{str},
    breaklines=true,
    numbers=left,
    numberstyle=\tiny\color{gray},
    language=C++,
    tabsize=4
}

% Seções numeradas independentemente (sem prefixo de capítulo)
\renewcommand{\thesection}{\arabic{section}}

% ---- Metadados ----
\autor{João Pedro dos Santos Henrique Plinta}
\titulo{Árvore B com Compressão de Dados}
\local{Pinhais}
\data{2026}
\orientador{Gabriel Vinicius Canzi Candido}
\curso{Bacharelado em Ciência da Computação}
\preambulo{Trabalho 2 apresentado à disciplina de Estruturas de Dados Avançadas do Curso de Bacharelado em Ciência da Computação do Instituto Federal do Paraná -- Campus Pinhais.}

% ------------------------------------------------------------------
\begin{document}

\pretextual
\imprimircapa
\tableofcontents
\cleardoublepage

\textual

% ------------------------------------------------------------------
\section{Introdução}
% ------------------------------------------------------------------

Este relatório descreve a implementação e os resultados do Trabalho 2 da disciplina de Estruturas de Dados Avançadas. O trabalho consiste em:

\begin{enumerate}
    \item Implementar uma \textbf{Árvore B} com chaves do tipo \textit{string}, com suporte a inserção, busca, remoção e persistência em arquivo binário;
    \item Implementar dois \textbf{algoritmos de compressão} --- LZW e Huffman Canônico --- para comprimir o arquivo gerado pela árvore;
    \item Medir e comparar o desempenho de ambas as estruturas em diferentes cenários.
\end{enumerate}

A motivação é simular o ambiente real de sistemas operacionais e bancos de dados, que utilizam árvores B para indexação eficiente e precisam persistir essas estruturas em dispositivos de armazenamento. A compressão reduz o espaço necessário e o tempo de leitura/escrita em disco.

O trabalho é composto por dois programas independentes escritos em \textbf{C++17}:

\begin{center}
\begin{tabular}{ll}
\toprule
Programa & Função \\
\midrule
\texttt{btree} & Gerencia a árvore B interativamente via linha de comando \\
\texttt{compress} & Comprime e descomprime o arquivo da árvore \\
\bottomrule
\end{tabular}
\end{center}

% ------------------------------------------------------------------
\section{A Árvore B}
% ------------------------------------------------------------------

\subsection{Definição e propriedades}

A Árvore B com grau mínimo $T$ satisfaz as seguintes propriedades:

\begin{itemize}
    \item Todo nó tem no máximo $2T - 1$ chaves;
    \item Todo nó interno (exceto a raiz) tem no mínimo $T - 1$ chaves;
    \item A raiz tem no mínimo 1 chave (se a árvore não for vazia);
    \item Todo nó interno com $k$ chaves tem exatamente $k + 1$ filhos;
    \item Todas as folhas estão no mesmo nível;
    \item As chaves em cada nó estão em ordem crescente.
\end{itemize}

A altura de uma Árvore B com $N$ chaves é $O(\log_T N)$, o que garante custo logarítmico para inserção, busca e remoção.

\subsection{Implementação}

A implementação utiliza as seguintes estruturas em C++:

\begin{lstlisting}
struct Node {
    int n;                        // numero de chaves
    bool leaf;
    std::vector<std::string> keys; // chaves[0..n-1]
    std::vector<Node*> ch;         // filhos[0..n]
    Node(int t, bool lf) : n(0), leaf(lf),
        keys(2*t-1), ch(2*t, nullptr) {}
};
\end{lstlisting}

As chaves são armazenadas como \texttt{std::string}, com comparação lexicográfica. O grau mínimo $T$ é definido na linha de comando.

\subsubsection{Inserção}

A inserção utiliza a estratégia \textit{proativa}: ao descer pela árvore, nós cheios ($n = 2T-1$) são divididos (\textit{split}) antes de continuar a descida. Isso evita que seja necessário subir novamente após a inserção na folha. A posição dentro de cada nó é encontrada com \texttt{std::lower\_bound} (busca binária). A complexidade é $O(\log T \cdot \log_T N) = O(\log N)$.

\subsubsection{Busca}

A busca usa busca binária (\texttt{std::lower\_bound}) para localizar a posição da chave dentro de cada nó. Ao encontrar a posição correta, desce para o filho correspondente ou retorna ``não encontrado'' em uma folha. Complexidade: $O(\log T \cdot \log_T N) = O(\log N)$.

\subsubsection{Remoção}

A remoção trata três casos:
\begin{enumerate}
    \item \textbf{Chave em folha}: remoção direta com deslocamento dos vizinhos;
    \item \textbf{Chave em nó interno com filho suficientemente cheio}: substitui pela chave predecessora ou sucessora e remove recursivamente;
    \item \textbf{Chave em nó interno com filhos ``magros''}: realiza fusão (\textit{merge}) dos filhos e desce recursivamente.
\end{enumerate}

\subsection{Persistência em arquivo}

O arquivo binário é escrito em \textbf{pré-ordem} (raiz, filho$_0$, filho$_1$, \ldots). Cada nó é serializado como:

\begin{center}
\begin{tabular}{lll}
\toprule
Campo & Tipo & Descrição \\
\midrule
\texttt{leaf} & \texttt{int} (4 B) & 1 se folha, 0 se interno \\
\texttt{n} & \texttt{int} (4 B) & número de chaves \\
\texttt{len$_i$} & \texttt{int} (4 B) & comprimento da $i$-ésima chave \\
\texttt{key$_i$} & \texttt{char[len$_i$]} & bytes da $i$-ésima chave \\
\bottomrule
\end{tabular}
\end{center}

O cabeçalho do arquivo armazena o grau $T$, permitindo carregar a árvore corretamente mesmo que o programa seja iniciado com um grau diferente. O uso de tamanho variável por chave evita o desperdício de bytes nulos que ocorreria com campos de tamanho fixo.

\subsection{Dois graus escolhidos}

Foram escolhidos dois graus bastante diferentes para observar o comportamento em dois extremos:

\begin{center}
\begin{tabular}{lcc}
\toprule
Propriedade & $T = 3$ & $T = 50$ \\
\midrule
Mín.\ de chaves por nó & 2 & 49 \\
Máx.\ de chaves por nó & 5 & 99 \\
Mín.\ de filhos por nó interno & 3 & 50 \\
Máx.\ de filhos por nó interno & 6 & 100 \\
Fator de ramificação & baixo & alto \\
Altura esperada (100k chaves) & $\approx 11$ & $\approx 3$ \\
\bottomrule
\end{tabular}
\end{center}

$T = 3$ representa uma árvore com nós pequenos e muitos níveis, similar a uma Árvore Binária de Busca balanceada (porém com fator de ramificação 3). $T = 50$ é o perfil típico de bancos de dados e sistemas de arquivos, onde cada nó corresponde a um bloco de disco e convém armazenar o máximo de chaves por acesso.

\subsection{Métricas coletadas}

Após cada operação de inserção, busca ou remoção, o programa exibe:

\begin{lstlisting}[language=bash, numbers=none]
[nos visitados: 6 | tempo: 1.28e-06 s | memoria: 4480 KB]
\end{lstlisting}

\begin{description}
    \item[Nós visitados] Incrementado no início de cada chamada recursiva que acessa um nó.
    \item[Tempo] Medido com \texttt{std::chrono::steady\_clock} antes e depois da operação.
    \item[Memória] Lida de \texttt{VmRSS} em \texttt{/proc/self/status} após a operação (memória física residente do processo).
\end{description}

% ------------------------------------------------------------------
\section{Algoritmos de Compressão}
% ------------------------------------------------------------------

\subsection{LZW (\textit{Lempel--Ziv--Welch})}

O LZW é um algoritmo de compressão baseado em dicionário sem parâmetros de configuração. O dicionário é inicializado com as 256 entradas correspondentes a cada byte possível e cresce dinamicamente ao longo da compressão.

\subsubsection{Compressão}

O algoritmo mantém uma \textit{string} corrente $w$ inicialmente vazia. Para cada byte $c$ lido:
\begin{itemize}
    \item Se $w + c$ está no dicionário, $w \leftarrow w + c$;
    \item Caso contrário, emite o código de $w$, adiciona $w + c$ ao dicionário e $w \leftarrow c$.
\end{itemize}
Ao final, emite o código de $w$.

A implementação usa códigos de \textbf{16 bits} (\texttt{uint16\_t}), permitindo até 65.536 entradas. Quando o dicionário está cheio, ele é congelado (sem \textit{reset}).

\subsubsection{Descompressão}

O descompressor mantém o mesmo dicionário (reconstruído localmente) e inverte o processo. O único caso especial ocorre quando o código recebido ainda não foi adicionado ao dicionário --- nesse caso, a entrada é $w + w[0]$ (onde $w$ é a última entrada processada).

\subsubsection{Formato do arquivo}

\begin{center}
\begin{tabular}{lll}
\toprule
Campo & Tamanho & Conteúdo \\
\midrule
Magic & 4 B & \texttt{LZW\textbackslash 0} \\
Tamanho original & 8 B & \texttt{uint64\_t} LE \\
Número de códigos & 8 B & \texttt{uint64\_t} LE \\
Códigos & $N \times 2$ B & \texttt{uint16\_t} LE cada \\
\bottomrule
\end{tabular}
\end{center}

\subsection{Huffman Canônico}

O algoritmo de Huffman atribui códigos de comprimento variável a cada símbolo (byte), com códigos mais curtos para símbolos mais frequentes. A variante \textbf{canônica} garante que, dados apenas os comprimentos dos códigos, os códigos podem ser reconstruídos deterministicamente --- eliminando a necessidade de serializar a árvore inteira no arquivo.

\subsubsection{Compressão}

\begin{enumerate}
    \item Contar a frequência de cada byte no arquivo de entrada;
    \item Construir a árvore de Huffman com uma fila de prioridade mínima;
    \item Coletar os comprimentos dos códigos para cada um dos 256 símbolos possíveis;
    \item Gerar os códigos canônicos: ordenar símbolos por (comprimento, valor) e atribuir códigos inteiros em ordem crescente, deslocando à esquerda a cada mudança de comprimento;
    \item Codificar os dados bit a bit (MSB primeiro dentro de cada byte).
\end{enumerate}

\subsubsection{Descompressão}

O descompressor lê os 256 comprimentos do cabeçalho, reconstrói os códigos canônicos e constrói uma \textbf{trie binária} estática em vetor. Cada nó da trie armazena dois índices de filho (bit~0 e bit~1) e o símbolo decodificado ($-1$ se nó interno). A decodificação percorre a trie bit a bit — ao atingir uma folha, emite o símbolo e retorna à raiz — eliminando a busca em mapa de \textit{strings} da abordagem ingênua e reduzindo o custo de $O(k \cdot \log n)$ para $O(k)$ por símbolo (onde $k$ é o comprimento do código).

\subsubsection{Formato do arquivo}

\begin{center}
\begin{tabular}{lll}
\toprule
Campo & Tamanho & Conteúdo \\
\midrule
Magic & 4 B & \texttt{HUF\textbackslash 0} \\
Tamanho original & 8 B & \texttt{uint64\_t} LE \\
Padding & 1 B & bits de preenchimento no último byte \\
Comprimentos & 256 B & 1 byte por símbolo (0 = ausente) \\
Dados & variável & fluxo de bits codificados \\
\bottomrule
\end{tabular}
\end{center}

O cabeçalho tem tamanho fixo de \textbf{269 bytes}, o que explica a menor eficiência em arquivos pequenos.

% ------------------------------------------------------------------
\section{Testes e Resultados}
% ------------------------------------------------------------------

\subsection{Ambiente de testes}

\begin{center}
\begin{tabular}{ll}
\toprule
Componente & Especificação \\
\midrule
Sistema operacional & Linux 7.0.0 \\
Compilador & g++, flags \texttt{-O2 -std=c++17} \\
Chaves ($N \leq 10.000$) & \textit{strings} aleatórias de 4--12 chars (letras + dígitos) \\
Chaves ($N > 10.000$) & sequenciais embaralhadas (\texttt{k000000000}--\texttt{k000999999}) \\
Semente & 42 (reprodutível) \\
\bottomrule
\end{tabular}
\end{center}

\subsection{Desempenho de inserção, busca e remoção}

Para cada combinação de $N$ e $T$, foram medidos os nós visitados e o tempo de todas as operações (inserções, buscas e remoções), com média calculada sobre o conjunto completo.

\begin{table}[H]
\centering
\caption{Desempenho médio por operação (média de insert + read + delete)}
\label{tab:insert}
\begin{tabular}{c rr rr r}
\toprule
 & \multicolumn{2}{c}{$T = 3$} & \multicolumn{2}{c}{$T = 50$} & Redução \\
\cmidrule(lr){2-3}\cmidrule(lr){4-5}
$N$ & Nós visit. & Tempo ($\mu$s) & Nós visit. & Tempo ($\mu$s) & em nós \\
\midrule
100        & 2,91 & 0,79 & 1,18 & 0,55 & $-$59\% \\
10.000     & 6,39 & 0,77 & 2,31 & 0,91 & $-$64\% \\
100.000    & 8,08 & 1,05 & 2,93 & 0,85 & $-$64\% \\
1.000.000  & 9,93 & 1,20 & 3,54 & 1,28 & $-$64\% \\
\bottomrule
\end{tabular}
\end{table}

O crescimento dos nós visitados segue a razão logarítmica esperada: ao aumentar $N$ de 100 para 1.000.000 (fator $10.000\times$), $T = 3$ passa de 2,91 para 9,93 nós ($\approx 3{,}4\times$) e $T = 50$ passa de 1,18 para 3,54 ($\approx 3\times$). A redução relativa de nós visitados se mantém estável em torno de $-$64\% para $N \geq 10.000$.

Ambos os graus usam busca binária dentro dos nós ($O(\log T)$ comparações por nó). O total de comparações de chave é $O(\log T \cdot \log_T N) = O(\log N)$, independentemente de $T$. Na prática, os tempos são próximos: para $N = 1.000.000$, $T = 3$ leva 1,20~$\mu$s e $T = 50$ leva 1,28~$\mu$s, pois $T = 50$ visita menos nós mas faz mais comparações por nó ($\lceil\log_2 99\rceil = 7$ vs $\lceil\log_2 5\rceil = 3$).

\subsection{Estrutura da árvore}

\begin{table}[H]
\centering
\caption{Estrutura da árvore após inserção de $N$ chaves}
\label{tab:stats}
\begin{tabular}{c rrc rrc}
\toprule
 & \multicolumn{3}{c}{$T = 3$} & \multicolumn{3}{c}{$T = 50$} \\
\cmidrule(lr){2-4}\cmidrule(lr){5-7}
$N$ & Altura & Total nós & Fill & Altura & Total nós & Fill \\
\midrule
100        & 3  &      31 & 58,1\% & 1 &      1 & 90,9\% \\
10.000     & 7  &   3.108 & 64,3\% & 3 &    145 & 69,6\% \\
100.000    & 9  &  31.210 & 64,1\% & 3 &  1.436 & 70,3\% \\
1.000.000  & 10 & 312.680 & 64,0\% & 4 & 14.575 & 69,3\% \\
\bottomrule
\end{tabular}
\end{table}

Com $N = 1.000.000$, $T = 50$ tem apenas 4 níveis e 14.575 nós, contra 10 níveis e 312.680 nós para $T = 3$ --- uma redução de $\approx 21\times$ no número de nós. O fill factor (fração de slots de chave ocupados) se estabiliza em $\approx 64\%$ para $T = 3$ e $\approx 70\%$ para $T = 50$ a partir de $N = 10.000$.

As médias da Tabela~\ref{tab:insert} incluem busca e remoção. O comportamento de busca e remoção é consistente com a inserção: $T = 50$ visita em média $\approx 64\%$ menos nós em todos os tamanhos. Buscas por chaves inexistentes percorrem a árvore até uma folha, visitando tantos nós quanto a altura. A remoção visita nós adicionais nos casos em que é necessário realizar \textit{merge} ou empréstimo de chaves entre nós irmãos.

\subsection{Memória utilizada}

\begin{table}[H]
\centering
\caption{Memória residente (VmRSS) ao final da inserção de $N$ chaves}
\label{tab:mem}
\begin{tabular}{crrc}
\toprule
$N$ & $T = 3$ (KB) & $T = 50$ (KB) & Redução \\
\midrule
100        &  4.232 &  4.232 & 0\%  \\
10.000     &  5.148 &  4.788 & $-$7\%  \\
100.000    & 13.492 &  9.928 & $-$26\% \\
1.000.000  & 97.012 & 62.068 & $-$36\% \\
\bottomrule
\end{tabular}
\end{table}

Com $T = 50$, a memória é \textbf{36\% menor} para 1.000.000 de chaves ($\approx 95$ MB vs $\approx 61$ MB). Embora cada nó com $T = 50$ seja individualmente maior (até 99 slots de chave), o número total de nós é $\approx 21\times$ menor, resultando em menor uso total de memória.

\subsection{Compressão}

As árvores foram comprimidas com ambos os algoritmos para $T = 3$ e $T = 50$.

\begin{table}[H]
\centering
\caption{Taxa de compressão — $T = 3$}
\label{tab:compress_t3}
\begin{tabular}{c r rr rr}
\toprule
 & & \multicolumn{2}{c}{LZW} & \multicolumn{2}{c}{Huffman} \\
\cmidrule(lr){3-4}\cmidrule(lr){5-6}
$N$ & Original & Comprimido & Taxa & Comprimido & Taxa \\
\midrule
100        &  1,3 KB    &  1,7 KB    & $0{,}81\times$ &  1,0 KB    & $1{,}29\times$ \\
10.000     &  141,6 KB  &  88,3 KB   & $1{,}60\times$ &  82,2 KB   & $1{,}72\times$ \\
100.000    & 1.610,9 KB & 384,8 KB   & $4{,}19\times$ & 637,8 KB   & $2{,}52\times$ \\
1.000.000  & 15,7 MB    &  4,4 MB    & $3{,}60\times$ &  6,5 MB    & $2{,}41\times$ \\
\bottomrule
\end{tabular}
\end{table}

\begin{table}[H]
\centering
\caption{Taxa de compressão — $T = 50$}
\label{tab:compress_t50}
\begin{tabular}{c r rr rr}
\toprule
 & & \multicolumn{2}{c}{LZW} & \multicolumn{2}{c}{Huffman} \\
\cmidrule(lr){3-4}\cmidrule(lr){5-6}
$N$ & Original & Comprimido & Taxa & Comprimido & Taxa \\
\midrule
100        &  1,1 KB    &  1,5 KB    & $0{,}72\times$ &  0,94 KB   & $1{,}17\times$ \\
10.000     &  118,4 KB  &  84,0 KB   & $1{,}41\times$ &  73,4 KB   & $1{,}61\times$ \\
100.000    & 1.378,3 KB & 353,3 KB   & $3{,}90\times$ & 548,2 KB   & $2{,}51\times$ \\
1.000.000  & 13,5 MB    &  3,9 MB    & $3{,}46\times$ &  5,7 MB    & $2{,}38\times$ \\
\bottomrule
\end{tabular}
\end{table}

\begin{table}[H]
\centering
\caption{Tempo de compressão e descompressão ($T = 3$, descompressão Huffman com trie)}
\label{tab:compress_time}
\begin{tabular}{c rr rr}
\toprule
 & \multicolumn{2}{c}{LZW} & \multicolumn{2}{c}{Huffman} \\
\cmidrule(lr){2-3}\cmidrule(lr){4-5}
$N$ & Compressão & Descompressão & Compressão & Descompressão \\
\midrule
100        &  0,60 ms &   0,11 ms &  0,14 ms &   0,16 ms \\
10.000     & 17,5 ms  &   3,53 ms &  3,64 ms &   2,09 ms \\
100.000    & 142 ms   &  10,0 ms  & 16,8 ms  &  22,3 ms  \\
1.000.000  & 1.535 ms &  73,8 ms  & 186 ms   & 210 ms    \\
\bottomrule
\end{tabular}
\end{table}

% ------------------------------------------------------------------
\section{Análise e Discussão}
% ------------------------------------------------------------------

\subsection{Impacto do grau $T$ no desempenho}

O resultado mais significativo é a diferença no número de nós visitados entre $T = 3$ e $T = 50$. Com 100.000 chaves, $T = 3$ visita em média \textbf{8,08 nós} por inserção, enquanto $T = 50$ visita apenas \textbf{2,93}. Isso reflete diretamente a diferença de altura:

\[
h_{T=3}(N) = \left\lceil \log_3 \frac{N+1}{2} \right\rceil
\qquad
h_{T=50}(N) = \left\lceil \log_{50} \frac{N+1}{2} \right\rceil
\]

Para $N = 100.000$: $h_{T=3} \approx 11$ e $h_{T=50} \approx 3$.

Em sistemas de banco de dados reais, onde cada nó representa um bloco de disco (tipicamente 4--16 KB), a redução de $\approx 8$ para $\approx 3$ acessos a disco por operação representa uma melhoria expressiva. Com latência de disco de $\approx 5$ ms por acesso, a economia seria de $\approx 25$ ms por operação.

No contexto em memória principal, ambos os graus usam busca binária dentro dos nós. Como o número total de comparações de chave é $O(\log N)$ independentemente de $T$, as diferenças de tempo são modestas. Para $N = 1.000.000$, $T = 50$ é ligeiramente mais lento (1,28~$\mu$s vs 1,20~$\mu$s): embora visite menos nós, cada nó com até 99 chaves exige $\lceil\log_2 99\rceil = 7$ comparações contra $\lceil\log_2 5\rceil = 3$ para $T = 3$, compensando parte do ganho em profundidade.

\subsection{Comparação entre LZW e Huffman}

O comportamento dos dois algoritmos varia conforme o tamanho do arquivo e o tipo de chave:

\begin{itemize}
    \item Para $N = 100$ e $N = 10.000$ (chaves aleatórias): Huffman supera o LZW em taxa ($1{,}29\times$/$1{,}72\times$ vs $0{,}81\times$/$1{,}60\times$ para $T=3$). O LZW chega a \textbf{expandir} o arquivo para $N = 100$ ($0{,}81\times$), pois o \textit{overhead} do cabeçalho supera o ganho em um arquivo muito pequeno;
    \item Para $N = 100.000$ e $N = 1.000.000$ (chaves sequenciais): LZW supera expressivamente o Huffman ($4{,}19\times$ vs $2{,}52\times$ para $T=3$, $N=100.000$). A alta regularidade das chaves sequenciais (\texttt{k000000000}, \texttt{k000000001}, \ldots) cria padrões de bytes repetitivos que o dicionário LZW explora muito eficientemente;
    \item \textbf{Queda do LZW em 1.000.000}: a taxa cai de $4{,}19\times$ ($N=100.000$) para $3{,}60\times$ ($N=1.000.000$). Com o arquivo maior, o dicionário de 65.536 entradas satura antes de processar todo o conteúdo e novos padrões não podem mais ser registrados, reduzindo a eficiência;
    \item \textbf{Tempo}: o Huffman é $\approx 10\times$ mais rápido na compressão para todos os tamanhos. A descompressão LZW é rápida ($\approx 74$~ms para 1M chaves); a descompressão Huffman leva $\approx 210$~ms via trie binária estática, que percorre bit a bit com custo $O(k)$ por símbolo (onde $k$ é o comprimento do código), sendo $\approx 7\times$ mais rápida que a abordagem por mapeamento de \textit{strings}.
\end{itemize}

\subsection{Integridade dos dados}

Em todos os testes, os arquivos descomprimidos foram verificados com \texttt{cmp} e são \textbf{bit a bit idênticos} aos originais. Isso confirma a corretude dos dois algoritmos de compressão e descompressão.

% ------------------------------------------------------------------
\section{Conclusão}
% ------------------------------------------------------------------

O trabalho demonstrou que a Árvore B é uma estrutura de indexação eficiente para até 1.000.000 de chaves, com operações na casa dos microssegundos e crescimento logarítmico no número de nós acessados.

A escolha do grau $T$ tem impacto direto na estrutura: com $T = 50$ e $N = 1.000.000$, a árvore tem apenas 4 níveis e 14.575 nós, contra 10 níveis e 312.680 nós para $T = 3$ ($21\times$ menos nós). Em memória principal, com busca binária dentro dos nós, o total de comparações é $O(\log N)$ para qualquer $T$. O ganho em nós visitados ($-64\%$) não se traduz em ganho proporcional de tempo: para $N = 1.000.000$, $T = 50$ é marginalmente mais lento (1,28 vs 1,20~$\mu$s), pois compensa a menor profundidade com mais comparações por nó. Em sistemas orientados a disco, onde cada acesso a nó representa um bloco físico, $T$ grande seria claramente superior.

Em relação à compressão, o desempenho depende fortemente do tamanho do arquivo e da regularidade das chaves. Para arquivos pequenos com chaves aleatórias ($N \leq 10.000$), o \textbf{Huffman} é superior em taxa e é $\approx 10\times$ mais rápido. Para arquivos grandes com chaves sequenciais ($N \geq 100.000$), o \textbf{LZW} supera o Huffman expressivamente em taxa ($4{,}19\times$ vs $2{,}52\times$ para $N = 100.000$), pois o dicionário acumula os padrões altamente repetitivos das chaves sequenciais. Porém, o LZW perde eficiência ao saturar seu dicionário de 65.536 entradas ($3{,}60\times$ para $N = 1.000.000$ vs $4{,}19\times$ para $N = 100.000$). O Huffman mantém tempo de compressão baixo e uso de memória estável em todos os cenários; sua descompressão via trie binária ($\approx 210$~ms para 1M chaves) é mais lenta que o LZW ($\approx 74$~ms) mas opera em tempo $O(k)$ por símbolo.

Como trabalho futuro, seria interessante avaliar: (1) LZW com códigos de largura variável para superar o limite de 65.536 entradas; (2) serialização em largura (\textit{level-order}) ao invés de pré-ordem, pois nós do mesmo nível têm conteúdo similar e poderiam ser melhor comprimidos juntos.

\end{document}
